% Eventual scalar receipts for the selected nonlinear step.
\section{Scalar prices for a nonlinear step}
\label{sec:ex-lattice-price-receipts}

This section proves the scalar scale estimates for the shipment, probe,
purchase, lodging, and credit fields of the finite-menu tariff. The allowance
and gain estimates come from a separate arrival receipt on the same selected
run; they are not consequences of these scalar inequalities. All limits below
are as the scale parameter $x$ tends to infinity, with the block length fixed.

Fix a block side $L$, an integer occupancy bound $N$, a positive strength
$\mu$, a positive relative weight $\lambda$, an exponent $\rho\in(-1,0)$,
and a fixed cost $C>0$. Write
\[
 b(x)=x+c_L,\qquad R(x)=\lceil(\log x)^2\rceil,\qquad
 M(x)=\exp\bigl((\log\log b(x))^2\bigr),
\]
where $c_L$ is fixed once $L$ is fixed. Set
\[
 f(x)=b(x)^{-1},\qquad y(x)=b(x)^{-2},
\]
\[
 k(x)=\bigl(N\mathfrak s(R(b(x)))+4\bigr)y(x),
 \qquad
 z(x)=2\exp\bigl(-\mu M(x)^2/16\bigr)
       +N\mathfrak s(R(x))y(x).
\]
Here
\[
 \mathfrak s(r)=e\,\bigl(2(16r+9)\bigr)^4,
 \qquad \mathcal R(x)=b(x)^\rho e^{\lambda C}.
\]

\begin{lemma}[Rate calculus and decaying costs]
\label{lem:ex-lattice-price-rate-calculus}
For a function $u$, say that it has rate $a$ if, for every $\delta>0$,
eventually $|u(x)|\leq b(x)^{a+\delta}$. Products add rates, sums of finitely
many terms keep the larger rate, and multiplication by a fixed constant does
not change a rate. Moreover, if $u$ has rate $a<0$, then $u(x)\to0$.

The functions $R(x)$ and every fixed polynomial in $R(x)$ have rate zero.
The same is true of $\mathfrak s(R(x))$ and $\mathfrak s(R(b(x)))$. If
$\mu>0$, then the first summand in $z(x)$ has rate less than
every real number; in particular it has rate $-2$. The fixed shift yields
\begin{equation}
 R(b(x))\leq 2R(x)\quad\text{eventually}.
 \label{eq:ex-lattice-pressure-radius}
\end{equation}
Consequently $k$ and $z$ have rate $-2$, as does
\[
\bigl(2R(x)+1\bigr)^4z(x).
 \]
\end{lemma}

\begin{proof}
For products, apply the two rate bounds with exponent slack $\delta/2$ and
use
\[b(x)^{a+\delta/2}b(x)^{c+\delta/2}=b(x)^{a+c+\delta}.\]
For sums,
use the same slack and the fact that $b(x)^{\delta/2}$ eventually exceeds
the fixed number of summands. Fixed constants are absorbed by any positive
slack. Since $b(x)\to\infty$, a negative rate implies convergence to zero by
choosing an intermediate exponent still below zero.

Since $b(x)=x+c_L$, the ratio $(\log b(x))^2/(\log x)^2$ tends to one.
Thus, eventually, $(\log b(x))^2\leq(3/2)(\log x)^2$ and
$(\log x)^2\geq2$. Therefore
\[
 R(b(x))\leq(\log b(x))^2+1
 \leq(3/2)(\log x)^2+1
 \leq2(\log x)^2\leq2R(x),
\]
as claimed. Also, the definition $R(x)=\lceil(\log x)^2\rceil$ and the
fixed-shift relation show that $R(x)$ is bounded by a constant times
$(\log b(x))^2$ for large $x$. Applying this same comparison with argument
$b(x)$ controls $R(b(x))$. Every fixed power of
$\log b(x)$ is smaller than $b(x)^\delta$ for every $\delta>0$; hence $R$
and its fixed polynomials have rate zero. The site factor is a fixed
quartic polynomial in its argument, so it also has rate zero, including
after the comparison \eqref{eq:ex-lattice-pressure-radius}.

Finally, $M(x)^2=\exp(2(\log\log b(x))^2)$. For every $A>0$,
\[
 \frac{\mu}{16}M(x)^2\geq A\log b(x)
\]
for all sufficiently large $x$, since with $t=\log b(x)$ the logarithm of
the ratio $M(x)^2/t$ is $2(\log t)^2-\log t\to\infty$. Taking $A$ as
large as desired proves the claimed decay of the Gaussian summand. The
remaining rate statements follow by the product and sum rules just proved.
\end{proof}

\begin{proposition}[Eventually purchased tariff]
\label{prop:ex-lattice-eventual-tariff}
Assume $\rho\in(-1,0)$, $\lambda>0$, and $\mu>0$. Then, for all
sufficiently large $x$, the following inequalities hold simultaneously:
\begin{align}
 0&\leq f(x), & 0&<y(x)\leq1, \label{eq:ex-lattice-price-basic}\\
 N\mathfrak s(R(x))e^{(\lambda+1)C}f(x)
   &\leq \mathcal R(x), \label{eq:ex-lattice-price-shipping}\\
 k(x)+e^{k(x)}k(x)&\leq f(x), \label{eq:ex-lattice-price-buy}\\
 (2R(x)+1)^4z(x)&\leq\tfrac12,
 &e^{z(x)}-1&\leq1.
 \label{eq:ex-lattice-price-lodging}
\end{align}
The lodging and credit cost is $z(x)$, whose Gaussian term is
\[2\exp(-\mu M(x)^2/16).\]
Here the fine monitor equals $M(x)$. Thus these are exactly the nonnegativity, probe, shipment, purchase,
lodging, and credit conditions required by the local tariff.
\end{proposition}

\begin{proof}
For $x$ large, $b(x)>1$, so $f=b^{-1}>0$ and $0<y=b^{-2}\leq1$.
The polynomial factor has rate zero, while $f$ has rate $-1$. Thus
\[
 N\mathfrak s(R(x))e^{(\lambda+1)C}f(x)
\]
has rate $-1$. Choose $q_0$ strictly between $-1$ and $\rho$. Eventually
the left side is at most $b(x)^{q_0}$, which is at most $b(x)^\rho$ for
large $x$. Multiplying the fixed factor $e^{\lambda C}$ on the right gives
\[
 N\mathfrak s(R(x))e^{(\lambda+1)C}f(x)
 \leq b(x)^\rho e^{\lambda C},
\]
as claimed.

By the rate-calculus lemma, $k$ has rate $-2$. In particular, eventually
$0\leq k\leq b^{-3/2}\leq1$. On a further tail,
$b^{1/2}\geq1+e$, and therefore
\[
 k+e k\leq(1+e)b^{-3/2}\leq b^{-1}.
\]
Since $e^k\leq e$ for $k\leq1$,
this proves \eqref{eq:ex-lattice-price-buy}. The product
$(2R+1)^4z$ has rate $-2$, so it is eventually at most $1/2$. Also
$z\to0$ and $z\geq0$; hence eventually $z\leq\log2$, which gives
$e^z-1\leq1$. This proves all the claims.
\end{proof}

\begin{remark}[Placement in the selected construction]
\label{rem:ex-lattice-price-receipt-placement}
On the selected run, the Gaussian pressure estimate supplies this tariff.
The arrival receipt supplies local activity, allowance, and gain bounds; the
reset receipt supplies separate rate payments for soft response. The delay
is enlarged to satisfy all finitely many transitions at each stage. The
finite-menu estimate uses arrival and reset as distinct inputs: neither is a
consequence of the tariff. Its shipment budget is the tariff inequality
proved above, combined with the predecessor's paid activity bounds.
\end{remark}
